\section{iGPT：Autoregressive Pretraining 能否迁移到 Pixels}

若 generative pretraining 的核心是对序列分布建模，那么图像也可被量化、展平并按像素或 code 顺序预测。问题在于图像具有二维局部结构，raster order 是人为选择；语言中合理的 causal factorization 未必是视觉最有效归纳偏置。

\subsection{把图像变成序列}

本节先把图像转成可由 decoder 处理的序列。iGPT 将图像表示为像素 token，训练 next-pixel prediction；模型既可生成 completion，也可把 hidden representation 用于 classification probing。

\lecturefigure{slide-22-image-gpt-question.jpg}{Image GPT 检验同一 autoregressive recipe 能否跨到图像。}{官方视频 23:42；Chen et al. 2020。}

\begin{knowledgebox}{读图：左边是问题，右边是 factorization}
图像被按固定顺序展开，模型预测下一个像素/离散值。这样可以直接复用 decoder-only Transformer，却把二维邻接转换成一维次序。成功说明 architecture 有通用性，效率和 sample quality 的差距则说明领域 bias 仍重要。
\end{knowledgebox}

\subsection{Completions：条件分布的视觉版本}

接下来用 completion 观察条件分布。给定图像上半部分，模型生成下半部分；多样输出表示存在多个 plausible continuation，而不是模型“恢复唯一真图”。读图前应先问输入约束保留了什么、哪些属性仍不确定，以及不同样本是否真正覆盖多种合理模式。

\lecturefigure{slide-23-igpt-completions.jpg}{iGPT completion 展示视觉条件生成的多模态不确定性。}{官方视频 25:31。}

\begin{knowledgebox}{读图：正确答案不是原图像素}
若上半张脸可对应多种下半部，逐像素 ground truth 只是一个样本。评价应考虑 perceptual plausibility、global consistency 和 diversity。只用 pixel MSE 会惩罚合理但不同的生成，也可能偏好模糊平均。
\end{knowledgebox}

\subsection{Feature learning}

进一步检查生成 objective 是否学到可用于分类的 feature，这是 transfer loop 的关键。Probe accuracy 随层或模型规模变化，可以显示 representation 在何处最可分。

\lecturefigure{slide-24-igpt-feature-learning.jpg}{iGPT 同时评估 generative samples 与 learned visual features。}{官方视频 26:24。}

\begin{knowledgebox}{读图：生成质量与表示质量是两条轴}
表格比较模型、层或 evaluation protocol。一个模型可生成好看图像却不提供最佳 linear probe，也可能相反。Representation learning 需要冻结 feature、统一 classifier 与 data split，避免把 probe capacity 当成 backbone quality。
\end{knowledgebox}

\subsection{本章小结}

iGPT 证明 autoregressive pretraining 可以迁移到 pixels，并同时产生 generation 与 representation。它也暴露 sequence order、视觉偏置和评价协议的重要性。

